\documentclass[11pt]{article}

\usepackage[margin=0.8in]{geometry}
\usepackage{graphicx}
\usepackage{listings}
\usepackage{xcolor}
\usepackage{hyperref}
\usepackage{float}
\usepackage{setspace}
\usepackage{booktabs}
\usepackage{array}

\hypersetup{
    colorlinks=true,
    linkcolor=blue,
    urlcolor=blue
}

\lstdefinestyle{code}{
    basicstyle=\ttfamily\small,
    breaklines=true,
    frame=single,
    columns=fullflexible,
    keepspaces=true,
    showstringspaces=false,
    backgroundcolor=\color{gray!8},
    rulecolor=\color{gray!40}
}

% Use this until a real screenshot is ready. Replace each placeholder with
% \includegraphics[width=...]{screenshots/your-file.png}.
\newcommand{\evidenceplaceholder}[1]{%
    \begin{figure}[H]
        \centering
        \fbox{\parbox[c][1.55in][c]{0.86\linewidth}{\centering\textit{Screenshot placeholder}\\[0.2in]#1}}
    \end{figure}%
}

\setlength{\parskip}{0pt}
\setlength{\parindent}{0.5in}

\title{CSE 565 Design of Experiments Project}
\author{Chase Coleman}
\date{}

\begin{document}
\begin{titlepage}
    \centering
    \doublespacing
    \vspace*{1in}
    {\Large\textbf{CSE 565 Design of Experiments Project}\par}
    \vspace{0.5in}
    Chase Coleman\par
    Arizona State University\par
    CSE 565: Software Verification and Validation\par
    Design of Experiments Project\par
    \today\par
\end{titlepage}

\setcounter{page}{2}
\doublespacing

\section{Introduction and Mobile Application Specification}

This project applies pairwise Design of Experiments (DOE) testing to a mobile application that collects five customer inputs. Pairwise testing reduces the full input space while ensuring that every possible pair of factor values occurs in at least one test case. The complete exhaustive space contains $5 \times 3 \times 4 \times 4 \times 5 = 1{,}200$ combinations.

\begin{table}[H]
\centering
\caption{Mobile application factors and levels.}
\begin{tabular}{p{1.35in}p{4.65in}}
\toprule
\textbf{Factor} & \textbf{Allowed levels} \\
\midrule
Type of Phone & iPhone 14; iPhone 13; Galaxy Z; Huawei Mate; Google Pixel 7 \\
Authentication & Fingerprint; Face recognition; Text Password \\
Connectivity & Wireless; 3G; 4G LTE; 5G Edge \\
Memory & 128 GB; 256 GB; 512 GB; 1 TB \\
Battery Level & $< 20\%$; 20--39\%; 40--59\%; 60--79\%; 80--100\% \\
\bottomrule
\end{tabular}
\end{table}

% TODO: Add a short explanation of why pairwise coverage is appropriate for this application.

\section{Task 1: Pairwise Test Cases Created by a DOE Tool}

\subsection{DOE Tool Selection}

I selected Microsoft PICT (Pairwise Independent Combinatorial Testing) as the DOE tool. PICT accepts a plain-text model of factors and levels and generates a compact pairwise test suite. I used PICT version 3.7.4 with its default pairwise strength (two-way coverage).

\subsection{PICT Model}

The factor model used by PICT is stored in \texttt{pict\_model.txt}. It contains the five factors and all permitted levels from the assignment specification.

\lstinputlisting[style=code]{pict_model.txt}

\subsection{PICT Generation and Results}

The following command generated the pairwise test suite:

\begin{lstlisting}[style=code]
pict pict_model.txt
\end{lstlisting}

PICT reported 175 required value-pairs and generated 26 test cases. The raw output is stored in \texttt{pict\_pairwise\_test\_cases.tsv}.

\evidenceplaceholder{Replace with a screenshot of the PICT command, model, and generated output.}

% TODO: Add a concise explanation of the generated suite and insert selected rows or the full table if required.

\section{Task 2: Pairwise Test Cases Created by a Generative AI Tool}

\subsection{Generative AI Tool Selection}

ChatGPT will be used to independently generate a second pairwise test suite. The prompt will provide the same factor/level specification used by PICT and will request a complete table containing only allowed values.

\subsection{Prompt Used}

\begin{lstlisting}[style=code]
Create a pairwise (2-way) test suite for a mobile application. Return a
complete tab-separated table with exactly these five columns: Type of Phone,
Authentication, Connectivity, Memory, and Battery Level.

Use only these levels:
- Type of Phone: iPhone 14, iPhone 13, Galaxy Z, Huawei Mate, Google Pixel 7
- Authentication: Fingerprint, Face recognition, Text Password
- Connectivity: Wireless, 3G, 4G LTE, 5G Edge
- Memory: 128 GB, 256 GB, 512 GB, 1 TB
- Battery Level: <20%, 20-39%, 40-59%, 60-79%, 80-100%

Every pair of values from every two different columns must occur in at least
one row. Do not use values outside these lists. Do not include duplicate rows.
After the table, state the number of test cases generated.
\end{lstlisting}

\subsection{Generated Test Cases and Evidence}

ChatGPT generated 27 test cases in tab-separated format. The exact prompt and
unedited response are preserved in \texttt{chatgpt\_prompt\_and\_response.md},
and the table used for validation is stored in
\texttt{chatgpt\_pairwise\_test\_cases.tsv}.

\evidenceplaceholder{Replace with a screenshot showing the exact ChatGPT prompt and generated test-case table.}

\section{Task 3: Assessment of the Two Test Suites}

Both suites will be independently validated rather than relying only on the generation tool's coverage claim. The verification checks that every test case is complete, every value is permitted by the specification, duplicate rows are identified, and every required pair is covered.

\begin{table}[H]
\centering
\caption{Pairwise-suite comparison. Complete this table after generating and verifying the AI suite.}
\begin{tabular}{p{2.5in}cc}
\toprule
\textbf{Measure} & \textbf{PICT suite} & \textbf{ChatGPT suite} \\
\midrule
Test cases generated & 26 & 27 \\
Invalid values & 0 & 0 \\
Duplicate rows & 0 & 0 \\
Required value-pairs & 175 & 175 \\
Pairs covered & 175/175 & 175/175 \\
Pairwise coverage & 100\% & 100\% \\
\bottomrule
\end{tabular}
\end{table}

Both suites were checked with \texttt{verify\_pairwise\_coverage.py}. The PICT suite contains 26 valid, non-duplicate test cases and covers all 175 required value-pairs. The ChatGPT suite contains 27 valid, non-duplicate test cases and also covers all 175 required value-pairs. The PICT suite is one test case smaller and is reproducible from the same model and command. The ChatGPT result achieved the same coverage, but required a precise prompt and independent verification.

The two suites therefore satisfy the pairwise DOE requirement. Their results show that generative AI can produce a valid pairwise suite for this small model, but its stated coverage should not be accepted without validation. A deterministic DOE tool provides a more repeatable generation workflow, while an AI tool can assist with quickly creating and explaining a candidate suite.

\evidenceplaceholder{Replace with a screenshot of the independent verification output for both suites.}

\section{Task 4: Assessment of the DOE Tool}

\subsection{Features and Functionality}

% TODO: Explain PICT's text-based model, pairwise generation, supported coverage strength, and command-line output.

\subsection{Scope Covered by the Tool}

% TODO: Explain which specification inputs PICT modeled and any assumptions or constraints not represented.

\subsection{Performance}

% TODO: Discuss the 26-case output, coverage result, generation time, and reduction from 1,200 exhaustive combinations.

\subsection{Ease of Use}

% TODO: Describe installation, model syntax, command-line execution, readability of output, and reproducibility.

\section{Task 5: Assessment of the Generative AI Tool}

\subsection{Prompting and Result-Processing Experience}

% TODO: Describe how the prompt was designed, whether revisions were required, and how the output was converted into a verifiable table.

\subsection{Significance of Generative AI for DOE Testing}

% TODO: Discuss the benefits of AI assistance and its limitations, especially the need for independent validity and coverage checks.

\section{Conclusion}

% TODO: Summarize the DOE approach, the results of both verified suites, the comparison, and the recommended tool for this task.

\section*{References}

\hangindent=0.5in Microsoft. (n.d.). \textit{PICT: Pairwise Independent Combinatorial Testing}. GitHub. Retrieved September 21, 2026, from \url{https://github.com/microsoft/pict}

% TODO: Add any course materials or other sources actually cited in the report.

\end{document}
